Due to space limitations, below we mention a subset of the most relevant
outputs from our team, grouped by area of expertise.

\paragraph{Statistical Neuroscience:} \textbf{MS} is Professor of Theoretical
Neuroscience and Machine Learning, and director of the Gatsby Computational
Neuroscience Unit at UCL.
%
A substantial component of his research focuses on the development of advanced
machine-learning based tools for neuroscience research.
%
His group published a series of new neuroinformatics tools designed to
characterise and understand population-scale activity using the large-scale
multielectrode recording methods being developed. These papers provided the
backbone for a new analytic approach that is now being employed and extended by
systems neuroscience laboratories worldwide (e.g., Yu et al., J Neurophysiol,
2009; Macke et al., Advanced State space \ldots, 2015).
%
This remains an active thread of his research (Duncker et al., PMLR, 2019 and
Soulat et al., NeuroIPS, 2021).
%
Many of the algorithms described in these papers form the basis of capabilities
to be incorporated into the pipeline described in
Section~\ref{sec:offlineAnalysisMethods}.

\textbf{JR} is a research engineer fellow at the GCNU.
%
He has ample experience using statistical method to model electrophysiological
time series. He has developed statistical methods to characterize the response
properties of visual cells to stimulation with natural images (Rapela et al.,
JOV, 2006; Rapela et al., Netw Comput Neural Syst, 2010), to understand
attention in humans using electroencephalographic recordings (Rapela et al.,
Neural Comput, 2018), and to investigate disfunctional population activity in
epilepsy using cortical Utah array recording in the human cortex (Rapela et
al., EBMC, 2019).
%
At the Gatsby Unit he is the main developer of the Sparse Variational Gaussian
Process Factor Analysis Python
package\footnote{\url{https://github.com/joacorapela/svGPFA}}, he contributes
statistical methods and visualisations to the AEON project (Campagner et al.,
CoSyNe, 2025) and, funded by
BBSRC\footnote{\url{https://gow.bbsrc.ukri.org/grants/AwardDetails.aspx?FundingReference=BB\%2FW019132\%2F1}}
he is leading the development of the
Bonsai.ML\footnote{\url{https://bonsai-rx.org/machinelearning/}} package, with
\textbf{GL} and \textbf{NG}.

\paragraph{NaLoDuCo experimentation:} \textbf{TB} is a Professor of
Neuroscience and Behaviour at the SWC, where he leads the Foraging Working
Group that develops the AEON platform.
%
Over the past 20 years, he has delivered groundbreaking insights into the
behavioural and circuit mechanisms underlying ethological behaviours.
%
He has extensively characterised neural circuits for instinctive escape
behaviours (Evans et al., Nature, 2018) and extended these
insights to spatial navigation (Shamash et al., Nature
Neuroscience, 2021; Shamash et al., Neuron, 2022).

\textbf{DC} is a GCNU/SWC joint research fellow in neuroscience. He has
integrated computational tools, neural recording, naturalistic behavioural
assays, and neural manipulation to explore active sensation and decision-making
(Campagner et al., eLife, 2016; J.  Neuroscience, 2019).
%
His recent work identified a crucial cortico-collicular circuit required for
orienting to shelters during escape (Campagner et al., Nature, 2022).
%
He leads the implementation of the AEON platform (Campagner et al., SfN, 2024;
Campagner et al., CoSyNe, 2025).

\textbf{CS} is a senior scientist at the AIND, where he is spearheading odour
learning NaLoDuCo experimentation. Throughout his career, he has pursued
questions at the interface of systems neuroscience, sensory physiology, and
synaptic integration.
%
He has investigated how sensory experience reshapes cortical representations
(Schoonover et al., Nature, 2021),
%
developed tools for long-term extracellular recordings in rodent cortex
(Schoonover et al., Nature, 2021), and
%
devised strategies to capture volitional, naturalistic mouse behavior in order
to study how animals explore and learn about their environment (Fink, Axel,
Schoonover, eLife 2019; Fink, Hogan, Schoonover, Neurobiol Learn Mem, 2024).

\paragraph{Large-scale data analysis and sharing:} \textbf{SV} an
associate director of data and outreach at the AIND and a leader in big and
open science.
%
She led the Visual Coding project in creating and analyzing the unprecedented
Allen Brain Observatory 2-photon dataset, recording the visual activity in over
63,000 neurons across 6 cortical areas, 4 cortical layers, and 14 transgenic
Cre lines that label specific neural populations (de Vries et al., Nature,
2020).
%
This dataset is a pioneering demonstration of the potential of open science,
and has been reused in over 60 publications (de Vries et al., eLife, 2023).
%
DV has led scientists in the analysis of this dataset (e.g., King et al.,
eNeuro, 2023), investigated the function of VIP neurons in visual processing
(Millman et al., eLife, 2020; sharing this dataset in DANDI), and studied the
relation between 2-photon imaging and electrophysiology (e.g., Siegle et al.,
eLife, 2021).

\textbf{JY} is a Senior Data Infrastructure Engineer at the AIND.
%
Prior to joining AIND, he developed robust data infrastructure supporting
cloud-based financial applications in the private sector.
%
Currently, at AIND, he serves as the lead architect responsible for designing
and implementing large-scale, cloud-based pipelines for processing behavioral
and neurophysiological data.

\paragraph{Spike sorting:} \textbf{AB} is an electrophysiology pipeline
development engineer at the AIND. He develops innovative tools and methodologies
for extracellular electrophysiology.
%
He is the co-creator and lead developer of SpikeInterface, a popular spike
sorting package (Buccino et al., eLife, 2020), ProbeInterface, a framework for
probe handling in electrophysiology (Garcia et al., Front Neuroinfor, 2022),
MEArec, a simulator of extracellular recordings (Buccino and Einvoll,
Neuroinformatics, 2021), and SpikeForest, a spike-sorting benchmarking package
(Magland et al., eLife, 2020).
%
In addition, his research has contributed to several aspects of
electrophysiology data analysis including the impact of spike collisions on
spike sorting performance (Garcia et al., eNeuro, 2022), motion estimation and
correction (Garcia et al., eNeuro, 2024; Watters et al., eNeuro, 2025), and
compression strategies for large-scale electrophysiology data (Buccino et al.,
J Neural Eng,2023).

\paragraph{Visualisations:} \textbf{DB} is a software engineer at the AIND,
where he develops scalable visualisation tools.
%
He co-developed the Data and Atlas websites for the International Brain
Laboratory, enabling intuitive visualization of brain-wide electrophysiological
data (Birman et al., bioRxiv, 2024), and he contributed to the development of
Pinpoint, a 3D web-based platform for planning multi-probe electrophysiology
experiments (Birman et al, eLife, 2023).

\paragraph{Experimental control:} \textbf{GL} is the director of NeuroGEARS
Ltd., and the creator of Bonsai, a reactive visual programming language widely
used for experimental control (Lopes et al., Front Neuroinform, 2015).
NeuroGEARS has strong links to the SWC, is the main developer of the AEON
platform, and is the business partner of the SWC/GCNU in the BBSRC funded
project mentioned above.
%
\textbf{NG} is a software engineer at the GCNU. He is highly experience in
Bonsai software development, authored of the BonZeb package (Guilbeault et al.,
Sci Rep, 2021) and is lead developer of the Bonsai.ML package mentioned above.
